\section{Theta kernel and normalization}
\label{app:theta}

This appendix supplies the theta facts used in the localization argument,
including all normalization factors.  For $t>0$ and $x\in\mathbb R$, define
\begin{equation}\label{eq:app-theta-definitions}
 \begin{gathered}
 \vartheta(t)=\sum_{n\in\mathbb Z}e^{-\pi n^2t},\qquad
 G(x)=e^{x/2}\vartheta(e^{2x}),\qquad
 \kappa(x)=\frac{G''(x)-G(x)/4}{2},\\
 C(x)=\cosh(x/2),\qquad W(x)=\frac{\kappa(x)}{2C(x)}.
 \end{gathered}
\end{equation}
We prove that $\kappa$ and $W$ are smooth, even and strictly positive, that
\begin{equation}\label{eq:app-theta-envelope}
 0<W(x)\le 40\exp\!\left(4|x|-\pi e^{2|x|}\right),
\end{equation}
and that
\begin{equation}\label{eq:app-theta-mass}
 \int_{\mathbb R}2C(x)\kappa(x)\,\mathrm dx=1.
\end{equation}
In particular, \eqref{eq:app-theta-mass} is equivalently
$\int_{\mathbb R}4C(x)^2W(x)\,\mathrm dx=1$; it is not a claim that
$W$ itself has unit integral.

\paragraph{Gaussian periodization.}
For completeness, the needed theta transformation follows directly from
the Gaussian.  In the Fourier convention
\[
 F_t(\xi)=\int_{\mathbb R}e^{-\pi t y^2}e^{-2\pi i\xi y}\,\mathrm dy,
\]
differentiation under the integral and integration by parts give
$F_t'(\xi)=-(2\pi\xi/t)F_t(\xi)$.  Also $F_t(0)=t^{-1/2}$: indeed,
the square of $\int_{\mathbb R}e^{-\pi y^2}\,\mathrm dy$ is, by polar
coordinates,
$2\pi\int_0^\infty r e^{-\pi r^2}\,\mathrm dr=1$.
Solving the differential equation therefore gives
\begin{equation}\label{eq:app-gaussian-fourier}
 F_t(\xi)=t^{-1/2}e^{-\pi\xi^2/t}.
\end{equation}
All differentiations and integrations by parts here are justified by
integrable polynomial multiples of a Gaussian.

The periodization
$P_t(y)=\sum_{n\in\mathbb Z}e^{-\pi t(y+n)^2}$ is smooth and
$1$-periodic, with locally uniformly convergent differentiated series.
Its $k$th Fourier coefficient, by an absolutely convergent change of
variables, is $F_t(k)$.  Thus
\[
 P_t(y)=t^{-1/2}\sum_{k\in\mathbb Z}
                  e^{-\pi k^2/t}e^{2\pi i k y}.
\]
For an explicit uniqueness justification, both sides are continuous and
have the same Fourier coefficients; convolution of their difference with
the Fej\'er kernel
$K_N(y)=(N+1)^{-1}|\sum_{r=0}^Ne^{2\pi i r y}|^2$ is zero.
These convolutions converge uniformly to the difference: $K_N\ge0$,
$\int_0^1K_N=1$, and on every closed set away from the integers,
$K_N=O((N+1)^{-1})$.  Uniform continuity supplies the remaining
neighborhood of the integers.  Evaluating at $y=0$ proves
\begin{equation}\label{eq:app-theta-inversion}
 \vartheta(t)=t^{-1/2}\vartheta(1/t).
\end{equation}
Termwise differentiation of $\vartheta$ is valid on every compact
subinterval of $(0,\infty)$.  Differentiating
\eqref{eq:app-theta-inversion} at $t=1$ yields
\begin{equation}\label{eq:app-theta-derivative-one}
 \vartheta'(1)=-\frac14\vartheta(1).
\end{equation}

\paragraph{Smoothness, parity and positivity.}
Equation~\eqref{eq:app-theta-inversion} gives $G(-x)=G(x)$, so $G$,
$\kappa$, and $W$ are smooth and even.  Put $t=e^{2x}$.  Direct
differentiation of \eqref{eq:app-theta-definitions} gives
\begin{equation}\label{eq:app-theta-kappa-derivative}
 \kappa(x)=e^{x/2}\bigl(2t^2\vartheta''(t)+3t\vartheta'(t)\bigr).
\end{equation}
The constant term of $\vartheta$ is annihilated.  Pairing its positive and
negative indices gives, in particular for $x\ge0$,
\begin{equation}\label{eq:app-theta-kappa-series}
 \kappa(x)=\sum_{j\ge1}
 \left(4\pi^2j^4e^{9x/2}-6\pi j^2e^{5x/2}\right)
             e^{-\pi j^2e^{2x}}.
\end{equation}
For $x\ge0$, each coefficient is
$e^{x/2}(4u^2-6u)>0$, where $u=\pi j^2e^{2x}\ge\pi>3$.
Consequently $\kappa>0$ on that half-line, and evenness gives positivity
everywhere.  Since $C>0$, the same holds for $W$.

\paragraph{A complete uniform envelope.}
The elementary bounds $3<\pi<22/7$ can be obtained without numerical
input from
\[
 \pi=4\int_0^1\frac{\mathrm du}{1+u^2}
 >4\sum_{k=0}^7\frac{(-1)^k}{2k+1}
 =\frac{135904}{45045}>3,
 \qquad
 \frac{22}{7}-\pi
 =\int_0^1\frac{u^4(1-u)^4}{1+u^2}\,\mathrm du>0.
\]
The first inequality follows by the finite geometric expansion with
positive remainder, and the second identity by polynomial division.
Also $e^3>\sum_{k=0}^4 3^k/k!=131/8>16$, so
$e^{-3\pi}<e^{-9}<1/4096$.

For $x\ge0$, use $2C(x)\ge e^{x/2}$ and discard the negative term in
\eqref{eq:app-theta-kappa-series}.  Since
$j^2-1\ge3(j-1)$, this gives
\begin{align*}
 W(x)
 &\le4\pi^2e^{4x-\pi e^{2x}}
           \sum_{j\ge1}j^4e^{-\pi(j^2-1)e^{2x}}\\
 &\le4\left(\frac{22}{7}\right)^2e^{4x-\pi e^{2x}}
           \sum_{j\ge1}\frac{j^4}{4096^{j-1}}.
\end{align*}
This bounds the full theta tail, not merely its first summand.
For every integer $j\ge1$, $j^4\le16^{j-1}$, by induction using
$((j+1)/j)^4\le16$.  Hence the coefficient above is bounded by
\[
 4\left(\frac{22}{7}\right)^2
   \sum_{j\ge1}\frac{j^4}{4096^{j-1}}
 \le4\left(\frac{22}{7}\right)^2\frac{256}{255}
 =\frac{495616}{12495}<40,
 \qquad 40-\frac{495616}{12495}=\frac{4184}{12495}>0.
\]
Evenness proves \eqref{eq:app-theta-envelope} on the whole line.  In
particular, $2C\kappa=4C^2W$ is integrable.

\paragraph{Exact mass.}
For $t\ge1$, the differentiated theta series gives
$\vartheta(t)\to1$ and $t^a\vartheta'(t)\to0$ for every fixed real $a$.
For example, after factoring out $e^{-\pi t}$, the remaining sums are
bounded by the convergent series
$\sum_{j\ge1}j^m e^{-\pi(j^2-1)}$, with a suitable fixed $m$.
Using evenness, $\mathrm dx=\mathrm dt/(2t)$, and
\eqref{eq:app-theta-kappa-derivative}, one obtains
\begin{align*}
 \int_{\mathbb R}2C\kappa\,\mathrm dx
 &=\int_1^\infty
    \left[2(t^{3/2}+t)\vartheta''(t)
                +3(\sqrt t+1)\vartheta'(t)\right]\,\mathrm dt\\
 &=\left[2(t^{3/2}+t)\vartheta'(t)+\vartheta(t)\right]_1^\infty\\
 &=1-4\vartheta'(1)-\vartheta(1)=1,
\end{align*}
where the last equality is exactly
\eqref{eq:app-theta-derivative-one}.  The displayed theta tail bounds
justify the endpoint limits and the integrations.  This proves
\eqref{eq:app-theta-mass} directly, with no Mellin transform or
spectral assumption.
